% Generated draft: requires mathematical and source review.
\documentclass[11pt]{article}
\usepackage[margin=25mm]{geometry}
\usepackage{fontspec}
\setmainfont{Latin Modern Roman}
\usepackage{amsmath,amssymb,mathtools,unicode-math}
\setmathfont{Latin Modern Math}
\usepackage{enumitem,needspace,xurl,hyperref}
\setlength{\parindent}{0pt}
\setlength{\parskip}{5pt}
\setlength{\emergencystretch}{3em}
\newcommand{\N}{\mathbb{N}}
\newcommand{\Z}{\mathbb{Z}}
\newcommand{\Q}{\mathbb{Q}}
\newcommand{\R}{\mathbb{R}}
\newcommand{\C}{\mathbb{C}}
\newcommand{\F}{\mathbb{F}}
\newcommand{\A}{\mathbb{A}}
\newcommand{\PP}{\mathbb{P}}
\newcommand{\E}{\mathbb{E}}
\newcommand{\eps}{\varepsilon}
\newcommand{\dd}{\,\mathrm{d}}
\begin{document}
\section*{UnsolvedMath ID: 2000015}
\textbf{Source problem code:} catalogue entry
\subsection{Definitions and mathematical statement}
The source asks to establish the status of elasticity theory relative to atomistic models. This is a foundational survey and derivation problem, not a single universally quantified theorem with supplied model hypotheses.

One precise kind of connection, mentioned by the supplied assessment, is a discrete-to-continuum $\Gamma$-limit. Let $\epsilon\downarrow0$ denote an atomic length scale, and suppose atomistic energies $E_\epsilon$ have been embedded or interpolated into a common topological deformation space $(\mathcal Y,\tau)$. To derive a continuum elastic energy $I$ by $\Gamma$-convergence is to prove both:
\begin{enumerate}
\item for every $y_\epsilon\to y$ in $\tau$, $I(y)\leq\liminf_{\epsilon\to0}E_\epsilon(y_\epsilon)$;
\item for every $y$, there exists $y_\epsilon\to y$ in $\tau$ with $\limsup_{\epsilon\to0}E_\epsilon(y_\epsilon)\leq I(y)$.
\end{enumerate}
Compactness/equicoercivity assumptions are additional ingredients when connecting minimizers of the discrete and continuum models. A local elastic candidate may take the form $I(y)=\int_\Omega W(\nabla y)\,\dd x$, but neither this form nor a particular microscopic interaction is prescribed by the source.

The requested status account must specify which atomic configurations, interaction laws, energy scalings, boundary conditions, deformation spaces, convergence topologies, and temperature regimes justify which continuum elastic conclusions, and where limitations remain. The supplied statement does not choose these data or restrict the task to static energetic $\Gamma$-limits; a single model theorem is therefore not silently substituted for its broader foundational question.
\subsection{Short English statement}
Assess how, and under what assumptions, continuum elasticity follows from atomistic descriptions. The supplied literature-survey assessment says nonlinear elastic discrete-to-continuum $\Gamma$-convergence is established for many crystalline cases, while full treatment of amorphous, defective and thermal settings remains open. It describes the subject as largely addressed with open aspects remaining, despite the registry's general open label. No supporting references or precise scope of those broad claims are supplied, so this draft records the survey question and its limitations rather than asserting a complete current state of the field.
\subsection{Sources}
\begin{itemize}
\item UnsolvedMath permanent ID 2000015, source code catalogue entry, ``Some Open Problems in Elasticity --- Atomistic foundations.'' Registry curation and metadata: Ulam AI and UnsolvedMath Contributors, CC BY 4.0; updated 2026-08-06.
\item Ball, source list ``problems on Elasticity'' (2002), Problem 15: \url{https://people.maths.ox.ac.uk/ball/Articles%20in%20Conference%20Proceedings%20and%20Books/JMB%202002%20re%20Marsden%2060th.pdf}. Supplied PDF-text extraction accessed 2026-07-29. The metadata says the source presents an open item and marks status/release and rights review NEEDS\_REVIEW.
\item The supplied LITERATURE-SURVEY assessment gives no specific crystalline-limit theorem, model citation, or reference for the claimed amorphous/defective/thermal gaps. No missing bibliography is invented.
\end{itemize}
\end{document}
